\documentclass[12pt]{article}
\usepackage[T1]{fontenc}
\usepackage[utf8]{inputenc}
\usepackage[margin=2.2cm]{geometry}
\usepackage{amsmath,booktabs,multirow,graphicx}
\usepackage[hidelinks]{hyperref}
\renewcommand{\thetable}{S\arabic{table}}
\renewcommand{\thesection}{S\arabic{section}}

\title{Supplementary material for\\ ``Choosing backup paths from hazard history: disaster-aware path selection for traffic-engineering recovery, tested on real cyclones and earthquakes''}
\author{Uday Pratap Singh, Bersha Kumari, Ebtasam Ahmad Siddiqui}
\date{}

\begin{document}
\maketitle

All tables below are generated by \texttt{paper/analysis.py} and \texttt{paper/supp\_tables.py} from result files committed in the repository \url{https://github.com/singhudaypratap/dart-disaster-aware-te}. The pre-registration trail was created in the development repository \url{https://github.com/singhudaypratap/nipa-book-chapter} (branch \texttt{claude/awesome-gauss-9xxnge}, directory \texttt{disaster-te}); commit identifiers below refer to that repository. The same history, with the original messages and timestamps, is mirrored in the code repository, where the file \texttt{PREREGISTRATION.md} maps each identifier below to its counterpart.

% ---------------------------------------------------------------------------------------------
\section{Development history}
\label{sec:s1}

The method was developed on Abilene (SNDlib capacities and demand) and BellSouth (Topology Zoo, uniform capacity, gravity demand). Training scenarios were twelve Atlantic hurricanes of 2004--2017 at 34 and 64~kt, 800 random disks per network and, for the variants marked ``+cuts'', all one- and two-link cuts. Evaluation used nine Atlantic hurricanes of 2018--2024 and 800 independently drawn disks per network. Table~\ref{tab:s1} lists every version that was run, with the values as they were seen at the time (raw, before the rounding-noise treatment of Section~\ref{sec:s3}); CVaR is given after that treatment.

\begin{description}
\item[Pilot (v1)] All $k-1$ backup tunnels chosen greedily by weighted coverage of training scenarios. The criterion (lower 95th percentile \emph{and} lower worst case than every classical rule on both networks, no-failure MLU within 5\%) was committed before the run (commit \texttt{88d30a8}) and \textbf{failed} on both networks (result committed verbatim, \texttt{ac90413}). Diagnosis: once the training footprints were covered, the remaining tunnels were chosen by tie-breaks, which overfitted the training events, and the coverage objective ignored capacity, which raised the BellSouth no-failure MLU from 0.60 to 0.76.
\item[v2] Second tunnel forced to be link-disjoint from the first; remaining tunnels by coverage (\texttt{29f2244}). Passed on BellSouth, failed on Abilene: one 50~km disk removed all of its closely routed backups.
\item[v3] Backbone of $k-1$ mutually link-disjoint tunnels; last tunnel by coverage (\texttt{0acf5de}). Same Abilene failure: on a network with node degree at most three, only three disjoint tunnels exist, and the data-free tie-break replaced the next-shortest path that the edge-disjoint rule keeps.
\item[v4 (frozen)] Edge-disjoint tunnels, last tunnel replaced only if a candidate strictly increases training coverage (\texttt{9da2e8f}). Lower mean and CVaR than every classical rule on both development networks; the single worst BellSouth disk remained higher (22.3\% against 19.0\% for edge-disjoint paths). Because one scenario is a fragile statistic, the tail criterion for the held-out tests was changed to $\mathrm{CVaR}_{0.95}$ before those tests were written. This change is a deviation from the pilot criterion and is reported here and in the main text.
\end{description}

\begin{table}[h]
\centering
\caption{Development history on the natural test scenarios of the development networks (Abilene: 6 hurricane and 800 disk scenarios; BellSouth: 16 and 800). Each version is shown with the edge-disjoint rule for reference.}
\label{tab:s1}
\resizebox{\linewidth}{!}{\begin{tabular}{lllrrrrr}
\toprule
Stage & Network & Rule & Mean (\%) & P95 (\%) & Worst (\%) & CVaR$_{0.95}$ (\%) & No-failure MLU \\
\midrule
Pilot (v1) & Abilene & DART v1 & 0.65 & 2.62 & 17.73 & 4.08 & 0.600 \\
Pilot (v1) & Abilene & Edge-disjoint & 0.05 & 0.54 & 3.04 & 0.66 & 0.600 \\
\addlinespace
Pilot (v1) & BellSouth & DART v1 & 0.66 & 4.35 & 12.95 & 6.72 & 0.758 \\
Pilot (v1) & BellSouth & Edge-disjoint & 0.46 & 1.85 & 18.96 & 8.22 & 0.600 \\
\addlinespace
Dev. v2 & Abilene & DART v2 & 0.04 & 0.00 & 17.73 & 0.80 & 0.600 \\
Dev. v2 & Abilene & Edge-disjoint & 0.05 & 0.54 & 3.04 & 0.66 & 0.600 \\
\addlinespace
Dev. v2 & BellSouth & DART v2 & 0.22 & 1.43 & 11.09 & 3.07 & 0.600 \\
Dev. v2 & BellSouth & Edge-disjoint & 0.46 & 1.85 & 18.96 & 8.22 & 0.600 \\
\addlinespace
Dev. v3 & Abilene & DART v3 & 0.04 & 0.00 & 17.57 & 0.78 & 0.600 \\
Dev. v3 & Abilene & Edge-disjoint & 0.05 & 0.54 & 3.04 & 0.66 & 0.600 \\
\addlinespace
Dev. v3 & BellSouth & DART v3 & 0.23 & 1.48 & 15.68 & 3.25 & 0.600 \\
Dev. v3 & BellSouth & Edge-disjoint & 0.46 & 1.85 & 18.96 & 8.22 & 0.600 \\
\addlinespace
Dev. v4 (frozen) & Abilene & DART v4 & 0.00 & 0.00 & 2.72 & 0.09 & 0.600 \\
Dev. v4 (frozen) & Abilene & Edge-disjoint & 0.05 & 0.54 & 3.04 & 0.66 & 0.600 \\
\addlinespace
Dev. v4 (frozen) & BellSouth & DART v4 & 0.19 & 1.00 & 22.30 & 3.38 & 0.600 \\
Dev. v4 (frozen) & BellSouth & Edge-disjoint & 0.46 & 1.85 & 18.96 & 8.22 & 0.600 \\
\bottomrule
\end{tabular}
}
\end{table}

% ---------------------------------------------------------------------------------------------
\section{Pre-registration records}
\label{sec:s2}

\begin{table}[h]
\centering
\caption{Commits that fix each criterion before the corresponding test, and the commits that record the outcome. Times are git commit times (UTC+0 as recorded by the repository); the corresponding pushes are visible in the GitHub history of the branch.}
\label{tab:s2}
\small
\begin{tabular}{p{4.6cm}p{2.0cm}p{7.5cm}}
\toprule
Test & Commit & Content \\
\midrule
Pilot (v1) criterion & \texttt{88d30a8} & \url{config/pilot.yaml}: criterion and train/test split \\
Pilot outcome: FAIL & \texttt{ac90413} & \url{reports/pilot_2026-11-08_FAIL.txt} \\
Test 1 criterion, method frozen & \texttt{0cfd320} & \url{config/test_dart_v4.yaml}; DART v4 code \\
Test 1 outcome: PASS & \texttt{bcc202f} & \url{reports/final_test_dart_v4_2026-11-08_PASS.txt} \\
Test 2 event-selection rule fix & \texttt{269fa4f} & missing-radii policy, fixed after counting events and before any evaluation \\
Test 2 criterion & \texttt{6594535} & \url{config/test_real_events.yaml} \\
Test 2 outcome: PASS & \texttt{d8f7e5e} & \url{reports/real_events_test_dart_v4_2026-11-08_PASS.txt} \\
\bottomrule
\end{tabular}
\end{table}

\paragraph{Test 1 criterion} Per network, on the natural test scenarios: (1) $\mathrm{CVaR}_{0.95}$ of routable loss of DART lower than that of every classical $k=4$ rule; (2) mean routable loss lower than every classical rule, with a paired two-sided Wilcoxon test against each rule significant at 0.05 after Holm correction over all nine comparisons; (3) no-failure MLU at most 1.05 times that of shortest paths. Pass required all three networks. Outcome: passed on AttMpls, BtNorthAmerica and G\'EANT.

\paragraph{Test 2 criterion} Pooled over seven networks, at 34~kt and MMI VII: negative mean paired difference and Holm-corrected two-sided Wilcoxon $p<0.05$ against each of the three classical rules. Secondary analyses (event-clustered, stricter thresholds, per network) were fixed in the same commit. Outcome: passed.

\paragraph{Change to event selection before Test 2} The first build of the Test 2 event set excluded any cyclone with a single synoptic track row lacking wind radii, which removed 633 of 721 western Pacific storms. Before any method was evaluated on real events, the rule was changed to drop only the rows without radii and never interpolate across gaps longer than 12~h (commit \texttt{269fa4f}). Only event counts had been seen at that point.

% ---------------------------------------------------------------------------------------------
\section{Linear-programming rounding noise}
\label{sec:s3}

The split ratios returned by the solver sum to one only to within about $10^{-16}$. Before the code fix in commit \texttt{d172304}, this produced reported losses of that size in 410 of the 1{,}335 real-event results; the smallest genuine loss was $3.1\times10^{-4}$ of demand. The pre-registered tests compared raw values, so such noise could turn ties into apparent wins or losses. Table~\ref{tab:s3} compares both treatments. Mean differences are identical, all comparisons remain significant, and the corrected counts show that the rules tie in most events and that DART rarely loses more. The main text reports the corrected values. The small difference between the raw counts here (e.g.\ 80/55/76) and in the archived Test~2 report (81/54/76) comes from re-reading the stored values with exact floating-point parsing; it concerns noise-level values only.

\begin{table}[h]
\centering
\caption{Test 2 primary comparisons with raw values (as pre-registered) and with values below $10^{-9}$ of demand set to zero. B/W/T: scenarios in which DART lost less, lost more, or tied.}
\label{tab:s3}
\begin{tabular}{lrrrrr}
\toprule
 & \multicolumn{2}{c}{Raw values (pre-registered)} & \multicolumn{2}{c}{Values $<10^{-9}$ set to 0} & \\
\cmidrule(lr){2-3}\cmidrule(lr){4-5}
DART vs & B/W/T & $p_{\mathrm{Holm}}$ & B/W/T & $p_{\mathrm{Holm}}$ & Mean diff. (pp) \\
\midrule
Shortest paths & 105/43/63 & 1.0e-11 & 67/1/143 & 1.0e-11 & -0.473 \\
Edge-disjoint & 80/55/76 & 2.6e-05 & 47/2/162 & 8.6e-08 & -0.134 \\
Geo-diverse & 77/60/74 & 3.8e-03 & 41/6/164 & 1.7e-06 & -0.368 \\
\bottomrule
\end{tabular}

\end{table}

% ---------------------------------------------------------------------------------------------
\section{Secondary results}
\label{sec:s4}

\begin{table}[h]
\centering
\caption{Test 2 secondary analyses. Event-clustered: each real event averaged over the networks it hit. Stricter thresholds: 64-kt winds and MMI VIII.}
\label{tab:s4b}
\resizebox{\linewidth}{!}{\begin{tabular}{lrlrrrr}
\toprule
Analysis & $n$ & DART vs & Mean diff. (pp) [95\% CI] & B/W/T & $p_{\mathrm{Holm}}$ & $r_{\mathrm{rb}}$ \\
\midrule
Event-clustered (34 kt / MMI VII) & 141 & Shortest paths & -0.333 [-0.473, -0.215] & 52/0/89 & 1.0e-09 & 1.00 \\
 &  & Edge-disjoint & -0.114 [-0.176, -0.064] & 36/0/105 & 3.3e-07 & 1.00 \\
 &  & Geo-diverse & -0.519 [-0.814, -0.271] & 38/4/99 & 5.8e-07 & 0.88 \\
\addlinespace
Stricter thresholds (64 kt / MMI VIII) & 56 & Shortest paths & -0.325 [-0.568, -0.125] & 13/0/43 & 4.2e-03 & 1.00 \\
 &  & Edge-disjoint & -0.051 [-0.112, -0.009] & 10/0/46 & 4.9e-03 & 1.00 \\
 &  & Geo-diverse & -0.281 [-0.513, -0.099] & 13/0/43 & 4.2e-03 & 1.00 \\
\bottomrule
\end{tabular}
}
\end{table}

\begin{table}[h]
\centering
\caption{Test 2 mean routable loss (\% of demand) per network at the primary thresholds.}
\label{tab:s4c}
\begin{tabular}{lrrrrr}
\toprule
Network & Shortest paths & Geo-diverse & Edge-disjoint & Shortest, $k=8$ & DART \\
\midrule
AttMpls & 0.913 & 0.487 & 0.417 & 0.583 & 0.312 \\
BtNorthAmerica & 0.833 & 0.056 & 0.274 & 0.246 & 0.070 \\
IIJ & 0.515 & 1.137 & 0.218 & 0.105 & 0.197 \\
CERNET & 0.493 & 0.082 & 0.327 & 0.080 & 0.003 \\
G\'EANT & 0.489 & 0.356 & 0.414 & 0.168 & 0.365 \\
GRNET & 0.000 & 0.000 & 0.000 & 0.000 & 0.000 \\
GARR & 0.000 & 0.000 & 0.000 & 0.000 & 0.000 \\
\bottomrule
\end{tabular}

\end{table}

\begin{table}[h]
\centering
\caption{All one- and two-link cuts on the Test 1 networks (1{,}596, 2{,}485 and 1{,}711 cuts): 95th percentile / worst routable loss (\% of demand), $k=4$.}
\label{tab:s4a}
\begin{tabular}{lrrrr}
\toprule
Network & Shortest paths & Geo-diverse & Edge-disjoint & DART \\
\midrule
AttMpls & 1.21 / 6.04 & 0.91 / 4.28 & 0.17 / 2.05 & 0.12 / 1.85 \\
BtNorthAmerica & 1.69 / 5.86 & 1.27 / 7.01 & 0.17 / 3.89 & 0.13 / 4.43 \\
G\'EANT & 3.46 / 23.34 & 1.57 / 14.35 & 1.06 / 12.30 & 0.77 / 9.83 \\
\bottomrule
\end{tabular}

\end{table}

\begin{table}[h]
\centering
\caption{Real test earthquakes (2018--2025) that failed at least one element. No rule lost routable traffic in any of them.}
\label{tab:s5}
\begin{tabular}{llrrr}
\toprule
Network & Scenario (USGS id, threshold) & Year & Failed nodes & Failed links \\
\midrule
BtNorthAmerica & \texttt{eq\_us70009gj8\_mmi7} & 2020 & 0 & 1 \\
G\'EANT & \texttt{eq\_us6000d3zh\_mmi7} & 2020 & 0 & 1 \\
G\'EANT & \texttt{eq\_us6000d3zh\_mmi8} & 2020 & 0 & 1 \\
G\'EANT & \texttt{eq\_us6000qznb\_mmi7} & 2025 & 0 & 2 \\
IIJ & \texttt{eq\_us70009gj8\_mmi7} & 2020 & 0 & 1 \\
IIJ & \texttt{eq\_us6000dher\_mmi7} & 2021 & 0 & 1 \\
GRNET & \texttt{eq\_us7000c7y0\_mmi7} & 2020 & 1 & 2 \\
GRNET & \texttt{eq\_us7000df40\_mmi7} & 2021 & 0 & 1 \\
\bottomrule
\end{tabular}

\end{table}

% ---------------------------------------------------------------------------------------------
\section{Reproduction}
\label{sec:s5}
From the repository root, with Python 3.12 or newer:
\begin{verbatim}
bash scripts/setup_server.sh config/stage2_real_disasters.yaml \
     config/test_real_events.yaml
.venv/bin/python -m dte.pilot --config config/test_dart_v4.yaml \
     --workers 32 --out results/t1
.venv/bin/python -m dte.realtest build \
     --config config/test_real_events.yaml --out results/t2
.venv/bin/python -m dte.realtest evaluate \
     --config config/test_real_events.yaml --out results/t2 --workers 32
PYTHONPATH=. .venv/bin/python paper/analysis.py
PYTHONPATH=. .venv/bin/python paper/supp_tables.py
PYTHONPATH=. .venv/bin/python paper/figures.py
\end{verbatim}
Downloaded ShakeMap versions are recorded with checksums; USGS can reissue ShakeMaps, so a later download may differ slightly from the versions used here.



\end{document}
